\chapter{Conclusion}
\label{ch:conclusion}

This thesis presented a multi-stage drone detection architecture in which the base detector is tuned for recall and confuser discrimination is handled by downstream learned stages. The design was validated component by component against comparison variants and ablations; the production stack (\S\ref{sec:production_stack}) integrates the winning component at each stage, with quantitative claims registered in the evidence ledger against the commands that produced them.

\section{Answers to the Research Questions}
\label{sec:rq_answers}

\paragraph{RQ1 (confuser-gap suppression).}
\emph{Yes.} In the classifier-comparison configuration (baseline RGB, v2 patch verifier; \S\ref{sec:cumulative}), the cascade reduces OOD confuser fire rate from 52.1\% (RGB detector alone) to 10.3\% with \texttt{scene\_aware\_v3more\_32feat} and to 0.8\% with the open-world \texttt{fusion\_no\_fn\_v1.1} (Table~\ref{tab:cum_confuser}). Hard-negative mining at the detector stage achieves at most $66.2\% \to 41.9\%$ on helicopters and is essentially inert on birds at small scale ($94.4\% \to 94.2\%$, \S\ref{sec:rgb_three_variants}). The cascade therefore catches what the detector cannot, without requiring the detector itself to be retrained against the bird category at unacceptable recall cost. The leakage-aware \texttt{robust6} classifier --- and its shipped grayscale-hardened extension \texttt{robust8} --- reproduces this suppression on the current \texttt{ft4}/\texttt{v3b} detectors. On an unseen confuser surface it cuts the classifier-only false-alert rate to $0.143$ and the best-composed cascade to $0.057$ (\S\ref{sec:feature_selection}, Table~\ref{tab:robust6_pipeline}), using six leakage-free features in place of \texttt{sa32}'s $32$.
% [source: ledger=cascade-confuser-collapse; eval=clfzoo_fnfn; run=eval/cumulative_halluc.py; cache=eval/results/_cumulative_halluc/confuser_fusion_no_fn_model_v1.1/summary.json; config=confuser_zoo_1280]

\paragraph{RQ2 (in-distribution cost and surface dependence).}
\emph{Surface-dependent, not architecture-intrinsic.} The comparison configuration that establishes this (baseline RGB, \texttt{sa32}, v2 patch verifier) shows the sign of the exchange flip with the surface. On Svanstr\"om paired data the cascade trades drone $F1$ for OOD safety: S1 $F1=0.950$ falls to S3 $F1=0.895$ at \texttt{patch\_thr}$=0.9$ (\S\ref{sec:patch_threshold}, Table~\ref{tab:patch_sweep}). On the real-video diagnostic the same cascade --- same classifier, same temporal smoother --- gains $F1$ over the Stage-1 RGB number ($0.760 \to 0.826$ for baseline RGB). It still cuts confuser FPR threefold (\S\ref{sec:realvideo_cascade}, Table~\ref{tab:cascade_segment}). The classifier choice is load-bearing: \texttt{fusion\_no\_fn\_v1.1} collapses real-video drone recall to $R=0.128$ while \texttt{scene\_aware\_v3more\_32feat} preserves it (\S\ref{sec:realvideo_classifier}). The exchange is therefore a joint property of the classifier and the evaluation surface, not of the cascade architecture. The production stack carries the same conclusion forward with the leakage-aware \texttt{robust8} classifier and the per-frame \texttt{mlp\_v5}/\texttt{mlp\_v5\_ir\_aligned} verifiers (\S\ref{sec:production_stack}); the full-cascade real-video re-run with those production verifiers in the alert-gate slot is a registered open item (Figure~\ref{fig:mlp_pipeline_placeholder}).
% [source: ledger=cascade-segment-recovers; eval=svan_s3_sa32_thr08; run=(EVIDENCE_LEDGER sweep); config=svan_iop_1280_s9]

\paragraph{RQ3 (cross-modal transfer).}
\emph{Partially --- the transfer holds where the discriminative content is silhouette-shaped, and degrades where it is not.} The IR detector, trained zero-shot on thermal data only, reaches drone $F1=0.636$ on grayscale-converted RGB video aggregate and $F1=0.837$ on the hardest bird-cluttered single clip (\S\ref{sec:grayscale}, Tables~\ref{tab:ir_grayscale}, \ref{tab:realvideo_seagull}). On that clip it ties the strongest RGB detector at $F1=0.840$ while cutting confuser FPPI by $3.2\times$ aggregate ($0.512 \to 0.158$). The transfer is preserved by the single-channel intensity input; the raw-3-channel-RGB control collapses to $F1=0.295$, isolating the grayscale conversion as the load-bearing step. The transfer is partial: on clean drone takeoff clips the dedicated RGB detectors win by 28--41~pp $F1$. The pattern is consistent with a detector that has learned silhouette structure strongly and visible-light texture weakly.
% [source: ledger=ir-grayscale-fallback; eval=vid_drone_ir_gray; run=eval/eval_video_tests.py; config=video_drone_iop]

\paragraph{Methodological thread (HITL co-development).}
The IR detector's precision / recall trajectory across six revisions on a fixed test split (Table~\ref{tab:ir_evolution}) characterises the dataset-state $\leftrightarrow$ model-state coupling. The V5 regression ($P: 0.895 \to 0.768$) is the case the protocol got wrong: a large batch of positives was injected in bulk instead of through the per-batch review loop, so the unreviewed batch carried unlabelled drones that V5 fires on as scored-FP. V6 and \texttt{Final} recover the precision by corpus cleanup. Reporting the regression is informative because it identifies a failure mode of HITL ingestion that is easy to repeat.
% [source: ledger=ir-version-progression; eval=ir_final_v5; run=eval/ir_version_comparison.py; cache=eval/results/ir_version_comparison/ir_comparison_test_640_2026-05-16T20-04-39.csv; config=ir_final_640]

\section{Production Stack and Future Work}
\label{sec:production_stack}


The methodological consequence that survives independently of any single number is that multi-modal drone-detection results are not comparable across studies without explicit scoring-rule disclosure; the 28~pp $F1$ swing the audit documents on Svanstr\"om for identical underlying detections is larger than the inter-system differences typically debated in the literature.

The production stack as of the time of writing is as follows. The RGB detector is the SelCom-plus-confuser fine-tune \texttt{Yolo26n\_selcom\_confuser\_ft4\_1280} (\texttt{ft4}), the only confuser-injection configuration that passed every drone-recall regression gate (\S\ref{sec:training_recipes}). The earlier RGB variants --- the \texttt{Yolo26n\_trained} baseline, \texttt{hardneg\_v3more}, \texttt{retrained\_v2}, and the \texttt{selcom\_mixed\_ft2\_1280}/\texttt{selcom\_960} fine-tunes --- are comparison points that motivated this pick (\S\ref{sec:rgb_three_variants}, \S\ref{sec:selcom}), not co-deployed models.

The IR branch is the \texttt{v3b} detector (\S\ref{sec:ir_evolution}), with a grayscale-RGB fallback path (\S\ref{sec:ir_grayscale_mode}) when no thermal camera is available, and as a parallel detection channel on bird-cluttered scenes even when thermal is available. Frame-level fusion uses the eight-feature leakage-aware trust classifier \texttt{robust8} (\S\ref{sec:feature_selection}) --- the six-feature \texttt{robust6} extended with \texttt{rgb\_mean\_conf}, an \texttt{is\_grayscale} flag, and a per-class \texttt{trust\_rgb} threshold ($\tau{=}0.20$) --- re-mined on the production \texttt{ft4}/\texttt{v3b} detectors and selected by the statistical leakage analysis; the $32$-feature \texttt{scene\_aware\_v3more\_32feat} (\texttt{sa32}), \texttt{control\_v3more\_40feat}, \texttt{fusion\_no\_fn\_v1.1}, and the base \texttt{robust6} are the comparison set against which \texttt{robust8} was validated (\S\ref{sec:classifier_compare}, \S\ref{sec:feature_selection}). The confuser verifier is the distilled \texttt{mlp\_v5} on the RGB branch and its cross-modal counterpart \texttt{mlp\_v5\_ir\_aligned} (one network, two per-modality scalers; \S\ref{sec:grayscale_verifier}) on the IR branch, both running per-frame in the deployed \texttt{MLPVerifier} and superseding the MobileNetV3 v2 patch verifier (\S\ref{sec:distill_verifier}) they beat on every confuser-rich surface and both modalities (\textsc{Ledger}~mlp-beats-patch-both-modalities). The two verifiers are composed trust-classifier-first (classifier$\to$verifier), the order that is most recall-safe at equal confuser suppression (\textsc{Ledger}~robust6-production-viable). The superseded patch verifier (at \texttt{patch\_thr=0.9} Svanstr\"om / $0.7$ real-video) remains the documented fallback on the photo-style \texttt{rgb\_dataset} surface where the MLP still trails (\S\ref{sec:distill_verifier}). The full-cascade real-video numbers in this thesis were measured with it, pending the \texttt{mlp\_v5} re-run (\S\ref{sec:realvideo_cascade}, Figure~\ref{fig:mlp_pipeline_placeholder}). In the production design the \texttt{mlp\_v5}/\texttt{mlp\_v5\_ir\_aligned} verifiers run \emph{per-frame} (always-on, reusing the detector's own ROI features; \textsc{Ledger}~v5-ship-per-frame); the \texttt{alert\_gate\_only} composition under which the real-video full-cascade numbers were measured is the as-measured \emph{patch}-verifier configuration noted above, pending the \texttt{mlp\_v5} re-run (Figure~\ref{fig:mlp_pipeline_placeholder}). Inputs are scored at \texttt{imgsz=1280} on Svanstr\"om-like content and \texttt{imgsz=960} on SelCom CCTV. Trust-aware scoring is the reporting rule.
% [source: ledger=robust8-grayscale-router,mlp-beats-patch-both-modalities,v5-ship-per-frame,ir-version-progression; eval=svan_gray_robust8_t20_clffilter; run=eval/compare_routing_pipeline.py; config=clf_own_holdout]

Three honest carve-outs accompany the stack: \texttt{robust6}'s IR and cross-modal features go to chance on the grayscale-fallback path (hence \texttt{robust8} adds \texttt{rgb\_mean\_conf} and a per-class \texttt{trust\_rgb} threshold at a measured real-video-confuser-fire cost; \textsc{Ledger}~robust8-grayscale-router, robust6-grayscale-ir-features-dead); \texttt{mlp\_v5} costs recall on the photo-style \texttt{rgb\_dataset} surface (\S\ref{sec:mlp_recall_drop}); and the aligned IR verifier over-vetoes on grayscale, so that path is gated rather than run unconditionally (\textsc{Ledger}~grayscale-trust-classifier-degrades). Latency on Jetson-class edge hardware is not yet measured; this and the real-video patch-threshold sweep are the open production-readiness items.
% [source: ledger=v5-ship-per-frame,robust8-grayscale-router,robust6-production-viable,mlp-beats-patch-both-modalities,ft4-backbone-freeze,ir-version-progression; eval=svan_gray_robust8_t20_clffilter; run=eval/overnight_ablation.py,eval/compare_routing_pipeline.py; cache=eval/results/_overnight_ablation/ablation_results.json; config=clf_own_holdout]

Four future-work directions: a track-based classifier consuming the temporal smoother's per-track state as a fourth evidence channel (the most promising single direction); closing the 2--7\% OOD-airplane gap with explicit Roboflow airplane crops as confuser-class data; a grayscale-aware routing gate for \texttt{mlp\_v5\_ir\_aligned}, which still over-vetoes on grayscale (\textsc{Ledger}~grayscale-trust-classifier-degrades); and edge-hardware latency benchmarking with TensorRT/ONNX quantisation. The open-world \texttt{fusion\_no\_fn\_v1.1} classifier remains selectable in place of \texttt{robust8} where confuser noise dominates and missed drones are operationally tolerable (\S\ref{sec:realvideo_classifier}).
% [source: ledger=patch-verifier-distribution-bound; eval=rob_rgb_baseline; run=eval/run_roboflow_eval.py; config=roboflow_rgb_drone_640]

% ======================================================================
% APPENDIX A — DATASETS
% ======================================================================
